% theory.tex -- Theory

\subsection{From Hopfield Energy to Score Function}
Let $\mathbf{X} = [\mathbf{m}_1, \dots, \mathbf{m}_K] \in \mathbb{R}^{d \times K}$ be the memory matrix whose columns are vectorized protein sequences from a family alignment of $K$ members, each represented in $d$ dimensions (after encoding and, optionally, PCA projection). The modern Hopfield energy~\cite{ramsauerHopfieldNetworksAll2021} is
\begin{equation}\label{eq:hopfield-energy}
    E(\boldsymbol{\xi}) = \tfrac{1}{2}\|\boldsymbol{\xi}\|^2 - \tfrac{1}{\beta}\log\sum_{k=1}^{K} \exp\!\bigl(\beta\,\mathbf{m}_k^\top \boldsymbol{\xi}\bigr),
\end{equation}
where $\beta > 0$ is an inverse temperature. The gradient is $\nabla E(\boldsymbol{\xi}) = \boldsymbol{\xi} - \mathbf{T}(\boldsymbol{\xi})$, where
\begin{equation}\label{eq:retrieval-map}
    \mathbf{T}(\boldsymbol{\xi}) := \mathbf{X}\,\operatorname{softmax}\!\bigl(\beta\,\mathbf{X}^\top\boldsymbol{\xi}\bigr)
\end{equation}
is the softmax attention retrieval map~\cite{ramsauerHopfieldNetworksAll2021}. The Boltzmann distribution $p_\beta(\boldsymbol{\xi}) \propto \exp(-\beta\,E(\boldsymbol{\xi}))$ has score function
\begin{equation}\label{eq:score}
    \nabla \log p_\beta(\boldsymbol{\xi}) = -\beta\,\nabla E(\boldsymbol{\xi}) = \beta\bigl(\mathbf{T}(\boldsymbol{\xi}) - \boldsymbol{\xi}\bigr).
\end{equation}
This score is \emph{exact} and computed by a single attention operation; no score network or training is required.

Applying the unadjusted Langevin algorithm (ULA)~\cite{robertsTweedieExponentialConvergence1996} to the potential $U = \beta E$ with step size $\tilde{\alpha} = \alpha/\beta$ yields the \emph{stochastic attention} update:
\begin{equation}\label{eq:sa-update}
    \boxed{\;
    \boldsymbol{\xi}_{t+1} = (1 - \alpha)\,\boldsymbol{\xi}_t + \alpha\,\mathbf{X}\,\operatorname{softmax}\!\bigl(\beta\,\mathbf{X}^\top\boldsymbol{\xi}_t\bigr) + \sqrt{\tfrac{2\alpha}{\beta}}\;\boldsymbol{\epsilon}_t, \quad \boldsymbol{\epsilon}_t \sim \mathcal{N}(\mathbf{0}, \mathbf{I})
    \;}
\end{equation}
Each step performs three operations: (i) contraction toward the origin, (ii) a softmax-weighted pull toward stored memories, and (iii) isotropic Gaussian noise scaled by $\sqrt{2\alpha/\beta}$. The per-step cost is $\mathcal{O}(dK)$: two matrix-vector products and one softmax. As $\beta \to \infty$ the noise vanishes and the update reduces to deterministic retrieval; as $\beta \to 0$ sampling becomes noise-dominated.

\subsection{Critical Temperature from PCA Dimensionality}\label{sec:beta-star}
The attention entropy $H(\beta) = -\sum_k p_k \log p_k$, where $p_k = \operatorname{softmax}(\beta\,\mathbf{X}^\top\boldsymbol{\xi})_k$, quantifies selectivity: $H = \log K$ when attention is uniform (generation) and $H \to 0$ when attention concentrates on a single memory (retrieval). The entropy derivative satisfies~\cite{alswaidanVarnerStochasticAttention2025}
\begin{equation}\label{eq:entropy-deriv}
    \frac{dH}{d\beta} = -\beta\,\mathrm{Var}_{\mathbf{p}}(e),
\end{equation}
where $e_k = \mathbf{m}_k^\top\boldsymbol{\xi}$ are the query-memory similarities. The maximum rate of entropy loss occurs at the inflection point $\beta^*$ satisfying $d^2H/d\beta^2 = 0$, which marks the boundary between the generation and retrieval regimes. The form of $\beta^*$ is constrained by the geometry of the stored patterns. Because the memory columns lie on the unit sphere $\mathbb{S}^{d-1}$, concentration of measure dictates that the similarity variance is $\mathrm{Var}(e_k) = 1/d$ exactly for any query drawn uniformly from $\mathbb{S}^{d-1}$ (Appendix~\ref{app:concentration}). The characteristic similarity scale is therefore $\sigma = 1/\sqrt{d}$, and since $\beta$ enters the softmax as $\beta e_k$, the inflection must occur when $\beta\sigma = \mathcal{O}(1)$, giving $\beta^* \sim \sqrt{d}$. A separate calculation shows that the softmax entropy over $K$ i.i.d.\ Gaussian scores undergoes its inflection at a dimensionless temperature $\tau^* \approx 1.6$ that is nearly independent of $K$ for $K \geq 30$ (Appendix~\ref{app:tau-star}), providing a universal baseline for the transition.

Empirically, we find a simple linear relationship across seven Pfam families and WW-domain subsamples spanning $K = 20$--$420$ and $d = 18$--$186$ (Section~5):
\begin{equation}\label{eq:beta-prediction}
    \beta^* \approx 1.57 + 0.28\sqrt{d},
\end{equation}
which captures 96\% of the variance ($R^2 = 0.962$). The intercept $1.57 \approx \tau^*$ recovers the universal softmax inflection baseline, confirming that the same dimensionless transition governs both idealized Gaussian scores and real protein similarities. The coefficient $0.28$ of $\sqrt{d} = 1/\sigma$ is substantially smaller than the pure Gaussian prediction $\tau^*/\sigma \approx 1.6\sqrt{d}$; this reduction arises because $\beta^*$ is computed at stored patterns $\mathbf{m}_k$, where the self-similarity $\mathbf{m}_k^\top\mathbf{m}_k = 1$ anchors one score at the maximum of $\mathbb{S}^{d-1}$ and makes concentration easier than for a random query (Appendix~\ref{app:stored-probes}). From a practical standpoint, Eq.~\eqref{eq:beta-prediction} enables fully automatic operation: given a seed alignment, $d$ is determined by PCA at 95\% variance, and $\beta^*$ follows without a temperature sweep.

\subsection{Regularity and Convergence}\label{sec:convergence}
The modern Hopfield energy has Lipschitz-continuous gradient with constant $L = 1 + \beta\sigma_{\max}^2/2$, where $\sigma_{\max} = \|\mathbf{X}\|_{\mathrm{op}}$, and satisfies the dissipativity condition $\langle \nabla E(\boldsymbol{\xi}), \boldsymbol{\xi} \rangle \geq \|\boldsymbol{\xi}\|^2/2 - M^2/2$ with $M = \max_k \|\mathbf{m}_k\|$~\cite{alswaidanVarnerStochasticAttention2025}. When $\beta\sigma_{\max}^2 < 2$, the energy is strictly convex and ULA iterates converge to $p_\beta$ in Wasserstein-2 distance at a geometric rate~\cite{dalalyanTheoreticalGuaranteesApproximate2017,durmusNonasymptoticAnalysisUnadjusted2017}.

In the protein setting, operating temperatures of interest satisfy $\beta\sigma_{\max}^2 \gg 2$, placing the sampler in the multimodal regime where the energy landscape has one local minimum per stored sequence. Following standard practice for multimodal Langevin sampling~\cite{bovier2004metastability}, we adopt a multi-chain protocol: $N_c$ independent chains are initialized near distinct stored patterns, exploiting fast intra-basin mixing while achieving inter-basin coverage through diverse initialization. The Metropolis-adjusted variant (MALA) provides an acceptance-rate diagnostic: acceptance rates near $1.0$ confirm that ULA discretization bias is negligible for the chosen step size $\alpha$.
